\section{Corrections to the record}
\label{sec:corrections}

This project's convention is that corrections are stated as corrections, in their own section, and
not silently patched into the artifacts they correct. The v0.5 study opened a twelve-entry register
against the v0.4 record; four of its entries never reached print, and this study adds five of its
own. All are listed here with the measurement that forces them.

\subsection{Entries carried forward from the v0.5 register that were never published}
\label{sec:corr-vfour}
\label{sec:corr-dialect}

\begin{description}

\item[C7 --- ``an undeclared half-suit scores nothing''.] The v0.4 paper justifies its second
  forcing horizon on the ground that an unclaimed half-suit is worth zero. It is not, in this
  harness: residue at the action cap is adjudicated by physical majority, and on a wrong
  post-horizon declaration the declaring team holds a mean \numResidueHeld{} of the six cards. The
  v0.4 paper contradicts itself on this point --- its own rules section states the dialect
  correctly. Neither paper printed the correction. It is printed here, and it matters beyond
  bookkeeping: the argument it supports is the one v0.5 removed as mechanism M8, so the mechanism
  was right for a reason its own justification got wrong.

\item[C9 --- restarting the willingness ladder.] The v0.4 paper argues that restarting the
  forced-endgame ladder lets early, safer declarations resolve cards and sharpen the later ones.
  The premise is never satisfied: essentially every forced endgame in this dialect contains exactly
  \numLadderLiveSets{} live half-suit (\numLadderOccasions{} occasions at seed \numLadderSeed{}), so
  there are no later declarations to sharpen.

\item[The fitted adversary that ``fails to reach 50\%''.] The v0.4 findings document reports that a
  policy fitted specifically against v0.4 fails to reach parity with it. The same study's own local
  best-response table records the response achieving $51.19\%$ with a confidence interval of
  $[49.67, 52.72]$, and the paper states the null result correctly. The findings document turns a
  null result into a win. This is the single most consequential unretracted claim in the corpus,
  because it is the only place anything resembling an exploitability number has ever been reported.

\item[``Modelling declaration timing as optimal stopping is worth measurable win rate''.] The same
  document. The same study's own ablation puts the value-based stopping rule at $+0.12$ points with
  an interval of $[-1.23, +1.47]$ against a fixed threshold, and the v0.4 specification says in
  terms that the ablation does not resolve a benefit.

\end{description}

\subsection{Corrections this study adds}
\label{sec:corr-self}

\begin{description}

\item[The tie channel is not a defect, and it is not where it was said to be.] The v0.5 design
  register, and the recon that opened this study, priced the exact-tie channel as the largest
  available loss and located it in the target dimension. Both halves are wrong.
  \vsixTieSameSuitFive\% of ties are two \emph{cards} of one half-suit at one target, not one card
  at two targets; and the exact posterior separates \vsixTieExactSepPctFive\% of them, so no
  tie-break rule can beat enumeration order --- measured, all four rules realise the same hit rate
  to three significant figures. The associated ``+3.16 half-suits per game of headroom'' is a
  hindsight bound and prices clairvoyance. See \S\ref{sec:diag-irreducible}.

\item[Enumeration order decides a fifth of contested decisions, not more than half.] \vfive{}'s
  top-$K$ chain and threat pass re-scores the leaders and moves the pick at
  \vsixTieShipMovedFive\% of ties. The frequently-quoted figure conflates the tie rate with the
  rate at which array order is the deciding rule.

\item[``Exact inference costs six points'' is a statement about the prior, not about exactness.]
  Scored as predictors on identical states, the exact posterior under a uniform-deal prior is the
  \emph{worst} of the three inference paths measured, at \vsixBeliefExactNLL{} nats and
  \vsixBeliefExactHit\% argmax hit rate against the deployed approximation's
  \vsixBeliefShippedNLL{} and \vsixBeliefShippedHit\%. The policy prior is the entire difference.
  The matched-budget refit under the exact belief that three studies have listed as the outstanding
  experiment is therefore not worth running: the object it would fit is a less informative
  posterior. See \S\ref{sec:belief-exactworse}.

\item[The v0.5 fit resolved nothing, and the shipped vector is v0.4's.] \S\ref{sec:diag-fitnoise}.
  The consequence for the reader of the v0.5 paper is that its fitting section describes a procedure
  that did not, at the budget and step size used, move the policy; and the consequence for this
  study is that a working optimiser had to be built before any mechanism could be evaluated.

\item[The exact reference engine was unsound in multi-agent play.] The block dynamic program parked
  every instance's tables in a shared per-thread pool, so any second construction on the same thread
  silently repointed the first instance's tables: \numAliasMismatch{} mismatches in
  \numAliasChecks{} checks. Every previously published number computed with the exact belief in a
  configuration where more than one seat holds such an object is affected. The deployed
  approximate-belief numbers are not, because that path never queries the object --- which is
  exactly why the defect survived three studies. See \S\ref{sec:engine-repairs}.

\end{description}

\subsection{A note on comparability}
\label{sec:corr-comparability}

Two of the corrections above change how earlier numbers should be read rather than what they are,
and one changes the numbers themselves. Readers comparing across versions should note that this
study's evidence standard is not the corpus's: every claim here is paired, and every claim that
survived is re-run on a second, disjoint seed bank at a per-cell budget three to seven times larger
than the published ablation rows of the v0.3, v0.4 and v0.5 studies. \S\ref{sec:protocol-paired}
gives the measurements that forced that change, including five mechanisms in this study that cleared
a $95\%$ interval at the corpus's budget and returned zero at this one.
